\SourceLeaf{128}
\begin{center}
\textbf{REVIEW.}
\end{center}

\medskip
\noindent\rule{\linewidth}{.3pt}

\noindent Allgemeine Literatur-Zeitung for 1808. Halle-Leipzig 1808. No.\,45. February 12. pp.\,353--358.

\noindent\rule{\linewidth}{.3pt}
\medskip

Dresden, at Walther's Court Bookshop: \textsc{Leonelli}'s \emph{logarithmische Supplemente, als ein Beitrag, Mängel der gewöhnlichen Logarithmentafeln zu ersetzen. Translated from the French with some additions by} \textsc{Gottfried Wilhelm Leonhardi}, \emph{second lieutenant in the Electoral Saxon Field Artillery Corps.} 1806. 88 pp. in octavo.

This little book contains two independent papers with a common purpose. Both are intended to accomplish something that ordinary logarithmic tables, without further aids, cannot achieve as accurately or as conveniently. The first is meant to use some readily adaptable auxiliary tables supplied here to calculate the logarithm of any given number without too much effort, to more decimal places than ordinary tables permit. The second develops the idea of a special table for finding, in a single operation, the logarithms of sums or differences of two quantities given only by their logarithms. It is meant to make the calculation more convenient than other procedures and, according to the author's plan, likewise to give an unusually large number of decimal places (14). The book includes only a sample of this table so far. Since numerical work with logarithms and anything that facilitates it are of great importance to anyone who does much numerical calculation, it should be worth giving these investigations more than a
